% =============================================================================
%  Ghi chú dựng báo cáo DevOps — goldcast
%
%  Biên dịch bằng XeLaTeX (Overleaf: Menu → Compiler → XeLaTeX).
%
%  Nguyên tắc của file này: chỉ ghi những gì có thật trong repo.
%    - Mọi đoạn code lấy nguyên văn từ file trong Git, kèm mã commit.
%    - Mốc thời gian lấy từ `git log` và lịch sử pull request trên GitHub.
%    - Số đo lấy từ docs/nhat-ky.md và docs/logs/.
%    - Chỗ nào repo chưa có số liệu thì để \canbosung{...} màu đỏ,
%      không điền số ước đoán.
%
%  Phạm vi: quy trình DevOps tuần 2–5 (tới commit 3037e69, 01/10/2026).
%  Không gồm phần mô hình học máy (PR #8) và phần giao diện (PR #7).
% =============================================================================
\documentclass[12pt,a4paper]{article}

\usepackage{fontspec}
\usepackage[vietnamese]{babel}
\setmainfont{TeX Gyre Termes}
\setmonofont{DejaVu Sans Mono}[Scale=0.85]

\usepackage[margin=2.5cm]{geometry}
\usepackage{xcolor}
\usepackage{listings}
\usepackage{booktabs}
\usepackage{tabularx}
\usepackage{graphicx}
\usepackage{hyperref}
\hypersetup{colorlinks=true, linkcolor=blue!50!black, urlcolor=blue!50!black}

% Chỗ còn thiếu số liệu — tìm "canbosung" để rà trước khi nộp.
\newcommand{\canbosung}[1]{\textcolor{red}{[CẦN BỔ SUNG: #1]}}
\newcommand{\commit}[1]{\texttt{#1}}

% ---------------------------------------------------------------- listings ---
\definecolor{codebg}{RGB}{248,248,246}
\definecolor{codekw}{RGB}{0,72,140}
\definecolor{codecm}{RGB}{110,110,110}
\definecolor{codestr}{RGB}{150,60,0}

\lstset{
  basicstyle=\ttfamily\small,
  backgroundcolor=\color{codebg},
  frame=single, rulecolor=\color{black!20},
  numbers=left, numberstyle=\tiny\color{black!45}, numbersep=6pt,
  breaklines=true, columns=fullflexible, keepspaces=true,
  showstringspaces=false, extendedchars=true,
  keywordstyle=\color{codekw}\bfseries,
  commentstyle=\color{codecm}\itshape,
  stringstyle=\color{codestr},
  captionpos=b, aboveskip=1em, belowskip=1em,
}

\lstdefinelanguage{Dockerfile}{
  morekeywords={FROM,AS,RUN,COPY,WORKDIR,USER,ENV,ARG,EXPOSE,ENTRYPOINT,CMD},
  sensitive=true, morecomment=[l]{\#}, morestring=[b]",
}
\lstdefinelanguage{YAML}{
  morekeywords={true,false},
  sensitive=false, morecomment=[l]{\#}, morestring=[b]", morestring=[b]',
}
\lstdefinelanguage{Jenkinsfile}{
  morekeywords={pipeline,agent,any,stages,stage,steps,dir,sh},
  sensitive=true, morecomment=[l]{//}, morestring=[b]', morestring=[b]",
}
\lstdefinelanguage{Env}{
  sensitive=true, morecomment=[l]{\#},
}

\title{Ghi chú báo cáo đồ án DevOps\\\large goldcast — theo dõi và dự báo giá vàng}
\author{\canbosung{Họ tên, MSSV, lớp}}
\date{Cập nhật theo repo tại commit \commit{3037e69} (01/10/2026)}

\begin{document}
\maketitle
\tableofcontents
\newpage

% =============================================================================
\section{Tổng quan hệ thống}
% =============================================================================

goldcast gồm ba dịch vụ, chạy bằng Docker Compose:

\begin{center}
\begin{tabularx}{\linewidth}{lXl}
\toprule
Dịch vụ & Công nghệ & Image nền \\
\midrule
\texttt{frontend} & Next.js 15 (App Router, TypeScript, Tailwind) & \texttt{node:22-alpine} \\
\texttt{backend}  & Spring Boot 3.3, Java 21, Flyway             & \texttt{eclipse-temurin:21-jre-alpine} \\
\texttt{db}       & PostgreSQL 16                                & \texttt{postgres:16-alpine} \\
\bottomrule
\end{tabularx}
\end{center}

Hạ tầng CI/CD đã dựng trong đồ án:

\begin{itemize}
  \item \textbf{GitHub Actions} — CI được quản lý: chạy test backend cho mỗi pull
        request và mỗi lần push vào \texttt{main}.
  \item \textbf{Branch protection} trên \texttt{main}: bắt buộc pull request và
        check \texttt{backend-test} xanh, không cho admin vượt rào.
  \item \textbf{Jenkins} tự vận hành trong Docker (cổng 8081): chạy test, rồi dựng
        Docker image backend và frontend.
\end{itemize}

% =============================================================================
\section{Dòng thời gian thực tế}
% =============================================================================

Lấy từ \texttt{git log main} và lịch sử pull request.

\begin{center}
\small
\begin{tabularx}{\linewidth}{llX}
\toprule
Thời điểm & Commit / PR & Nội dung \\
\midrule
21/09/2026 & \commit{cae1e86} & Commit khởi tạo mã nguồn \\
23/09/2026 & \commit{bb4d99e} & Thêm nhật ký và khung CI \\
24/09/2026 & \commit{66820fe} & Workflow GitHub Actions chạy test backend \\
24/09/2026 & \commit{dcb3040} & Thêm test cố tình sai để kiểm chứng cổng CI \\
24/09/2026 & \commit{9561da2} & Gỡ test cố tình sai \\
24/09/2026 & \commit{20e414b} & Thử push thẳng vào \texttt{main} (kiểm chứng branch protection) \\
25/09/2026 & \commit{8ab21c5} & Thêm \texttt{Jenkinsfile} (stage test) \\
25/09/2026 & \commit{584d72a} & Image Jenkins có Maven \\
27/09/2026 & PR \#1, \#2, \#3 & Merge Jenkinsfile, nhật ký tuần 3, nhật ký tuần 4 \\
01/10/2026 & PR \#4 (\commit{a89dd0b}) & Jenkins dựng được Docker image backend \\
01/10/2026 & PR \#5 (\commit{606555f}) & Thêm stage dựng image frontend \\
01/10/2026 & PR \#6 (\commit{3ebd18a}) & Thêm buildx để Jenkins dùng BuildKit \\
02/10/2026 & PR \#9 & Nhật ký tuần 5 \\
\bottomrule
\end{tabularx}
\end{center}

% =============================================================================
\section{Tuần 2 — Chạy goldcast trên máy cá nhân}
% =============================================================================

\subsection{Cấu hình tách khỏi mã nguồn (12-Factor: Config)}

Cấu hình không viết cứng trong Java. Nếu viết cứng, mật khẩu bị ``hàn'' vào
image khi build: đọc ra rất dễ, nhưng muốn đổi thì phải build lại.

Đường đi của một biến cấu hình gồm ba chặng:

\begin{lstlisting}[language=Env, caption={Chặng 1 — \texttt{.env} (không commit, đã có trong \texttt{.gitignore})}]
APP_INGEST_SEED_ON_EMPTY=true
\end{lstlisting}

\begin{lstlisting}[language=YAML, caption={Chặng 2 — \texttt{docker-compose.yml} đọc \texttt{.env} và truyền vào container}]
  backend:
    environment:
      APP_INGEST_SEED_ON_EMPTY: ${APP_INGEST_SEED_ON_EMPTY:-true}
\end{lstlisting}

\begin{lstlisting}[language=YAML, caption={Chặng 3 — \texttt{application.yml} của Spring đọc biến môi trường}]
app:
  ingest:
    seed-on-empty: ${APP_INGEST_SEED_ON_EMPTY:true}
\end{lstlisting}

\subsection{Docker Compose}

\begin{lstlisting}[language=YAML, caption={\texttt{docker-compose.yml} (rút gọn: dịch vụ \texttt{backend})}]
  backend:
    build:
      context: ./backend
      dockerfile: Dockerfile
    container_name: goldcast-backend
    depends_on:
      db:
        condition: service_healthy
    environment:
      SPRING_DATASOURCE_URL: jdbc:postgresql://db:5432/${POSTGRES_DB:-goldcast}
      SPRING_DATASOURCE_USERNAME: ${POSTGRES_USER:-goldcast}
      SPRING_DATASOURCE_PASSWORD: ${POSTGRES_PASSWORD:-goldcast}
      APP_INGEST_ENABLED: ${APP_INGEST_ENABLED:-true}
      APP_INGEST_SEED_ON_EMPTY: ${APP_INGEST_SEED_ON_EMPTY:-true}
      APP_CORS_ALLOWED_ORIGINS: ${APP_CORS_ALLOWED_ORIGINS:-http://localhost:3000}
      TZ: Asia/Ho_Chi_Minh
    ports:
      - "${BACKEND_PORT:-8080}:8080"
    healthcheck:
      test: ["CMD", "wget", "-qO-", "http://localhost:8080/actuator/health"]
      interval: 10s
      timeout: 5s
      retries: 20
      start_period: 40s
    restart: unless-stopped
\end{lstlisting}

Hai điểm đáng trình bày:
\begin{itemize}
  \item \texttt{depends\_on} với \texttt{condition: service\_healthy}: backend chỉ
        khởi động khi Postgres đã sẵn sàng, chứ không chỉ khi container đã chạy.
        Frontend cũng chờ healthcheck của backend theo cùng cách.
  \item Cú pháp \verb|${BIEN:-mac_dinh}|: chạy được ngay cả khi thiếu \texttt{.env}.
\end{itemize}

\subsection{Dockerfile backend — multi-stage build}

\begin{lstlisting}[language=Dockerfile, caption={\texttt{backend/Dockerfile}}]
# ---- build ----
FROM maven:3.9-eclipse-temurin-21 AS build
WORKDIR /build

# Dependencies resolve in their own layer, so editing source does not re-download Maven Central.
COPY pom.xml .
RUN mvn -B -q dependency:go-offline

COPY src ./src
RUN mvn -B clean package -DskipTests

# ---- run ----
FROM eclipse-temurin:21-jre-alpine
WORKDIR /app

# tzdata so Asia/Ho_Chi_Minh resolves; the alpine image ships without it.
RUN apk add --no-cache tzdata \
    && addgroup -S app \
    && adduser -S -G app app

COPY --from=build /build/target/*.jar app.jar
RUN chown -R app:app /app
USER app

EXPOSE 8080

ENTRYPOINT ["sh", "-c", "exec java -XX:MaxRAMPercentage=75 -XX:+UseSerialGC -Djava.security.egd=file:/dev/./urandom -jar /app/app.jar"]
\end{lstlisting}

Các điểm kỹ thuật cần giải thích:
\begin{itemize}
  \item \textbf{Hai stage}: image chạy chỉ chứa JRE và file \texttt{.jar}, không
        mang theo Maven hay mã nguồn.
  \item \textbf{Tách lớp phụ thuộc}: \texttt{COPY pom.xml} và
        \texttt{dependency:go-offline} đứng trước \texttt{COPY src}, nên sửa code
        không làm tải lại thư viện.
  \item \textbf{Chạy bằng user thường} (\texttt{USER app}), không chạy bằng root.
  \item \texttt{-XX:MaxRAMPercentage=75}: JVM tính heap theo giới hạn bộ nhớ của
        container, không theo RAM của máy chủ.
\end{itemize}

\subsection{Số đo tuần 2}

Lấy từ dòng \texttt{time} cuối mỗi file trong \texttt{docs/logs/}.

\begin{center}
\begin{tabular}{llr}
\toprule
File log & Lệnh & Thời gian \\
\midrule
\texttt{tuan02-build-nguoi.txt} & \texttt{docker compose build --no-cache} & 2 ph 12 s \\
\texttt{tuan02-build-nong.txt}  & \texttt{docker compose up -d --build} & 36,0 s \\
\texttt{tuan02-build-lan1.txt}  & \texttt{docker compose up -d --build} & 19,5 s \\
\texttt{tuan02-khoi-dong.txt}   & \texttt{docker compose up -d} (không build) & 17,1 s \\
\bottomrule
\end{tabular}
\end{center}

\canbosung{Log \texttt{tuan02-build-lan1.txt} không có output của bước build, chỉ có
dòng ``Built''. Rất có thể image đã có sẵn từ trước nên con số 19,5 s không phải
thời gian build lần đầu. Cần xác nhận lại trước khi đưa vào báo cáo; ô ``Thời gian
build lần đầu'' trong nhật ký tuần 2 hiện cũng đang trống.}

% =============================================================================
\section{Tuần 3 — Continuous Integration với GitHub Actions}
% =============================================================================

\subsection{Workflow}

\begin{lstlisting}[language=YAML, caption={\texttt{.github/workflows/ci.yml} (commit \commit{66820fe})}]
name: CI

on:
  push:
    branches: [main]
  pull_request:
    branches: [main]

jobs:
  backend-test:
    runs-on: ubuntu-latest
    steps:
      - uses: actions/checkout@v4
      - uses: actions/setup-java@v4
        with:
          java-version: '21'
          distribution: 'temurin'
          cache: maven
      - name: Chay test backend
        working-directory: backend
        run: mvn -B test
\end{lstlisting}

\subsection{Kiểm chứng cổng chất lượng}

\begin{enumerate}
  \item Thêm \texttt{CiGateTest} chứa \texttt{assertEquals(1, 2)}
        (commit \commit{dcb3040}) $\rightarrow$ CI \#4 đỏ.
  \item Gỡ test đó (commit \commit{9561da2}) $\rightarrow$ CI \#5 xanh.
  \item Bật branch protection cho \texttt{main}, rồi push thẳng
        (commit \commit{20e414b}) $\rightarrow$ bị từ chối với
        \texttt{protected branch hook declined}.
\end{enumerate}

Ảnh minh chứng ghi trong nhật ký: \texttt{tuan03-ci-do.png},
\texttt{tuan03-branch-protection.png}.
\canbosung{hai ảnh này chưa có trong \texttt{docs/anh/} của repo}

\subsection{Số đo — ảnh hưởng của cache Maven}

Cùng commit \commit{20e414b}, chạy lại 4 lần bằng ``Re-run all jobs''.

\begin{center}
\begin{tabular}{lllr}
\toprule
Lần chạy & Cache & Job \texttt{backend-test} & Total duration \\
\midrule
Nguội & trống  & 34 s & 1 ph 15 s \\
Nóng 1 & có sẵn & 21 s & 26 s \\
Nóng 2 & có sẵn & 19 s & 24 s \\
Nóng 3 & có sẵn & 23 s & 29 s \\
\bottomrule
\end{tabular}
\end{center}

Kết luận: cache giảm thời gian job từ 34 s xuống 19–23 s (khoảng 38\%).
Dùng cột \emph{Job} chứ không dùng \emph{Total duration}, vì ở lần nguội có 41
trên 75 giây là thời gian chờ GitHub cấp máy ảo.
Hạn chế: lần nguội chỉ có một mẫu.
Dung lượng cache: \canbosung{số MB, xem tab Caches của repo trên GitHub}.

\subsection{Vấn đề gặp phải}
\begin{enumerate}
  \item Branch protection không có hiệu lực trên repo private gói Free
        $\rightarrow$ chuyển repo sang public, sau khi rà lịch sử Git để chắc
        chắn \texttt{.env} và khoá SSH chưa từng bị commit.
  \item Admin vẫn push thẳng được $\rightarrow$ bật
        ``Do not allow bypassing the above settings''.
  \item Đọc nhầm Total duration ở lần đo đầu (kết luận sai là nhanh 3,4 lần).
\end{enumerate}

% =============================================================================
\section{Tuần 4 — Jenkins, CI tự vận hành}
% =============================================================================

\subsection{Image Jenkins có Maven}

Image \texttt{jenkins/jenkins} chính thức không có Maven. Thay vì cài bằng giao
diện, môi trường CI được mô tả bằng Dockerfile để nằm trong Git và dựng lại được.

\begin{lstlisting}[language=Dockerfile, caption={\texttt{jenkins/Dockerfile} ở tuần 4 (commit \commit{584d72a})}]
FROM maven:3.9-eclipse-temurin-21 AS maven

FROM jenkins/jenkins:lts-jdk21

USER root
COPY --from=maven /usr/share/maven /opt/maven
RUN ln -s /opt/maven/bin/mvn /usr/local/bin/mvn
USER jenkins
\end{lstlisting}

\subsection{Pipeline đầu tiên}

\begin{lstlisting}[language=Jenkinsfile, caption={\texttt{Jenkinsfile} ở tuần 4 (commit \commit{8ab21c5})}]
pipeline {
    agent any

    stages {
        stage('Test backend') {
            steps {
                dir('backend') {
                    sh 'mvn -B test'
                }
            }
        }
    }
}
\end{lstlisting}

Job \texttt{goldcast-ci} kiểu Pipeline, đọc pipeline từ SCM, nhánh
\texttt{*/main}, trigger Poll SCM với lịch \texttt{H/5 * * * *}.

\subsection{Số đo}

\begin{center}
\begin{tabular}{lllr}
\toprule
Lần chạy & Trạng thái & Build duration & Maven \texttt{Total time} \\
\midrule
\#2 (nguội) & SUCCESS & 46 s  & \canbosung{số giây} \\
\#3 (nóng)  & SUCCESS & 5,2 s & 1,693 s \\
\bottomrule
\end{tabular}
\end{center}

Kiểm chứng tự kích hoạt: sau khi merge PR, build \#4 tự chạy trong vòng 5 phút,
ghi \emph{Started by an SCM change}, thời gian 6,1 s.

\subsection{So sánh hai mô hình CI}

\begin{center}
\begin{tabularx}{\linewidth}{lXX}
\toprule
& GitHub Actions & Jenkins \\
\midrule
Hạ tầng & Do GitHub quản lý & Tự vận hành (container trên máy cá nhân) \\
Workspace & Máy ảo mới mỗi lần & Dùng lại giữa các lần build \\
Lần chạy nóng & 19–23 s (vẫn phải biên dịch) & 5,2 s (Maven 1,7 s) \\
Công cụ & Runner có sẵn hàng chục công cụ & Phải tự thêm Maven, Docker \\
Kích hoạt & Webhook sự kiện GitHub & Poll SCM 5 phút/lần \\
Rủi ro & Chậm hơn & Build có thể xanh nhờ tàn dư của lần trước \\
\bottomrule
\end{tabularx}
\end{center}

\subsection{Vấn đề gặp phải}
\begin{enumerate}
  \item Năm plugin Pipeline cài hỏng chỉ vì một plugin gốc
        (\texttt{pipeline-groovy-lib}) bị timeout khi tải.
  \item Build đầu chết sau 1,6 s vì \texttt{main} chưa có \texttt{Jenkinsfile}:
        \texttt{git branch -d} không kiểm tra nhánh đã vào \texttt{main} chưa.
  \item Ô \texttt{base} của pull request mặc định không phải \texttt{main}.
  \item \texttt{README.md} bị xoá nhầm; CI vẫn xanh. Phát hiện nhờ đọc diff
        (21 dòng thêm, 265 dòng xoá), khôi phục ở commit \commit{1e92507}.
  \item Yêu cầu review chặn merge với đồ án một người $\rightarrow$ tắt review,
        giữ yêu cầu CI xanh.
  \item Poll SCM được tích nhưng ô Schedule trống nên không có build nào chạy.
\end{enumerate}

% =============================================================================
\section{Tuần 5 — Pipeline sinh ra artifact}
% =============================================================================

\subsection{Image Jenkins điều khiển được Docker}

\begin{lstlisting}[language=Dockerfile, caption={\texttt{jenkins/Dockerfile} hiện tại (commit \commit{3ebd18a})}]
FROM maven:3.9-eclipse-temurin-21 AS maven
FROM docker:27-cli AS docker

FROM jenkins/jenkins:lts-jdk21

USER root
COPY --from=maven /usr/share/maven /opt/maven
COPY --from=docker /usr/local/bin/docker /usr/local/bin/docker
COPY --from=docker /usr/local/libexec/docker/cli-plugins /usr/local/libexec/docker/cli-plugins
RUN ln -s /opt/maven/bin/mvn /usr/local/bin/mvn
USER jenkins
\end{lstlisting}

So với tuần 4 có thêm hai dòng \texttt{COPY}:
\begin{itemize}
  \item \texttt{/usr/local/bin/docker} (commit \commit{a89dd0b}): lệnh
        \texttt{docker} cho Jenkins. Container Jenkins được gắn
        \texttt{/var/run/docker.sock} để điều khiển Docker daemon của máy chủ.
  \item \texttt{cli-plugins} (commit \commit{3ebd18a}): plugin buildx, để
        \texttt{docker build} dùng BuildKit thay vì builder cũ.
\end{itemize}

\canbosung{lệnh \texttt{docker run} dựng container Jenkins (cổng 8081, volume
\texttt{jenkins\_home}, \texttt{-v /var/run/docker.sock:...}, \texttt{--group-add 0})
không nằm trong repo — nên chép lại vào báo cáo, hoặc đưa vào repo dưới dạng
script hay file compose}

\subsection{Pipeline hiện tại}

\begin{lstlisting}[language=Jenkinsfile, caption={\texttt{Jenkinsfile} hiện tại}]
pipeline {
    agent any

    stages {
        stage('Test backend') {
            steps {
                dir('backend') {
                    sh 'mvn -B test'
                }
            }
        }

        stage('Build backend image') {
            steps {
                sh 'docker build -t goldcast-backend:${BUILD_NUMBER} -t goldcast-backend:latest backend/'
            }
        }

        stage('Build frontend image') {
            steps {
                sh 'docker build -t goldcast-frontend:${BUILD_NUMBER} -t goldcast-frontend:latest frontend/'
            }
        }
    }
}
\end{lstlisting}

\begin{itemize}
  \item Stage dựng image đặt \emph{sau} stage test: không đóng gói thứ chưa qua
        kiểm thử. Đây là ranh giới giữa CI và CD.
  \item Mỗi image mang hai tag: số build (truy vết được) và \texttt{latest}.
\end{itemize}

\subsection{Số đo}

\begin{center}
\begin{tabular}{lllll}
\toprule
Build & Pipeline & Builder & Trạng thái cache & Thời gian \\
\midrule
\#5 & test + backend            & legacy   & --                         & 2 ph 35 s \\
\#6 & test + backend + frontend & legacy   & tải mới image nền frontend & 4 ph 05 s \\
\#7 & test + backend + frontend & BuildKit & cache BuildKit trống       & 1 ph 16 s \\
\#8 & test + backend + frontend & BuildKit & ấm hoàn toàn               & 6,5 s \\
\bottomrule
\end{tabular}
\end{center}

Ở trạng thái ổn định, toàn bộ pipeline (93 unit test và hai lần dựng image) mất
6,5 s. Bảng này \emph{không} cho phép kết luận BuildKit nhanh hơn builder cũ bao
nhiêu lần, vì giữa \#6 và \#7 có hai biến cùng đổi: builder và trạng thái cache.

\subsection{Điều tra tính tái lập của build}

\begin{center}
\begin{tabularx}{\linewidth}{llX}
\toprule
Cặp so sánh & Khác biệt thật & Nguyên nhân \\
\midrule
Tuần 2 với \#5 & có & Máy cá nhân dùng BuildKit, Jenkins dùng builder cũ \\
\#5 với \#6 & có & Nhãn \texttt{node:22-alpine} đã trỏ sang bản mới \\
\#7 với \#8 & \textbf{không} & Nội dung image giống hệt \\
\bottomrule
\end{tabularx}
\end{center}

Kiểm bằng \texttt{docker image inspect} cho cặp \#7/\#8: danh sách mã băm các
lớp, trường \texttt{Created}, \texttt{Config}, \texttt{History}, \texttt{RootFS}
giống hệt; chỉ khác \texttt{LastTagTime} và \texttt{Metadata.Ref}, là thông tin
của phiên build chứ không phải nội dung image.

Kết luận: nội dung image tái lập được nếu cố định công cụ build và image nền.
Artifact vẫn phải được giữ và luân chuyển thay vì build lại ở từng môi trường,
vì image nền và builder có thể thay đổi mà mã nguồn không đổi.

\subsection{Hạn chế đã nhận diện, chưa xử lý}
\begin{enumerate}
  \item \textbf{Image tích tụ}: mỗi build sinh một image mới và không tự xoá.
  \item \textbf{\texttt{NEXT\_PUBLIC\_API\_BASE\_URL} bị cố định trong image frontend}:
        stage dựng image không truyền \texttt{--build-arg}, nên image mang giá
        trị mặc định \texttt{http://localhost:8080}.
  \item \textbf{Pipeline không tự dựng lại được chính nó}: thay image Jenkins
        vẫn phải làm thủ công.
  \item \textbf{Quyền của Jenkins}: gắn \texttt{docker.sock} tương đương quyền
        root trên máy chủ.
\end{enumerate}

\subsection{Bài học phương pháp}

Bốn lần kết luận sai vì so sánh nhầm đại lượng: Total duration của GitHub
Actions (tuần 3), Total time của Maven (tuần 4), so \#6 với \#7 khi cache khác
nhau (tuần 5), và cột \texttt{IMAGE ID} (tuần 5). Cả bốn lần đều được phát hiện
bằng cùng một cách: khi con số không khớp dự đoán, mở ra xem đại lượng đó thật
sự đo cái gì.

% =============================================================================
\section{Danh sách cần bổ sung trước khi nộp}
% =============================================================================

\begin{itemize}
  \item Họ tên, MSSV, lớp ở trang bìa.
  \item Thời gian build lần đầu ở tuần 2 (xác nhận lại ý nghĩa của log
        \texttt{tuan02-build-lan1.txt}).
  \item Dung lượng cache Maven trên GitHub Actions (tuần 3).
  \item Maven \texttt{Total time} của Jenkins build \#2 (tuần 4).
  \item Ảnh minh chứng tuần 3–5: hiện \texttt{docs/anh/} chỉ có thư mục
        \texttt{Tuan2}. Khi có ảnh, chèn bằng:
\begin{lstlisting}[language={[LaTeX]TeX}]
\begin{figure}[h]
  \centering
  \includegraphics[width=0.9\linewidth]{anh/tuan03-ci-do.png}
  \caption{CI báo đỏ khi có test cố tình sai}
\end{figure}
\end{lstlisting}
  \item Lệnh \texttt{docker run} dựng container Jenkins.
\end{itemize}

\end{document}
